\section{Cài đặt môi trường và hướng dẫn chạy}
\label{sec:installation}
\subsection{Yêu cầu môi trường}
\begin{table}[H]\centering
\caption{Điều kiện để cài đặt ứng dụng}\label{tab:environment}
\begin{tabularx}{\linewidth}{@{}P{3.4cm}Y@{}}\toprule
Thành phần & Yêu cầu theo cấu hình dự án\\\midrule
Node / package manager & Node 22.12+ trong dòng 22; pnpm theo \code{.mise.toml} và lockfile.\\
PHP & PHP 8.3; extensions fileinfo, mbstring, PDO/PDO MySQL; Composer cho autoload.\\
Database & MySQL/InnoDB, utf8mb4; database ứng dụng và database riêng \code{_test}.\\
Web server & PHP dev server khi phát triển; Apache/Nginx cấu hình PHP khi triển khai.\\
Quyền filesystem & PHP ghi được uploads và session cache; không public cấu hình/secret.\\
LaTeX & TeX Live với LuaLaTeX, latexmk, Biber; Python 3 nếu dùng build script báo cáo.\\
\bottomrule\end{tabularx}
\end{table}
Ghi lại phiên bản PHP/MySQL thực tế khi cài; bản MySQL trên máy khác cần được kiểm tương thích với migrations và tests.

\subsection{Chuẩn bị source và cấu hình}
Thực hiện tại root \code{E-commerceWebsitePrototype}. Cài dependency theo lockfile, sau đó chuẩn bị PHP bằng script nếu máy chưa có toolchain:
\begin{lstlisting}[caption={Chuẩn bị môi trường phát triển trên PowerShell}]
pnpm.cmd install --frozen-lockfile
npm.cmd run backend:setup
Copy-Item backend/.env.example backend/.env
# Edit backend/.env with local configuration.
\end{lstlisting}
Chỉ sao chép \code{backend/.env} khi chưa tồn tại; không ghi đè cấu hình đang dùng. Setup tải PHP/Composer vào .tools; nếu đã có PHP thì cài autoload bằng Composer và kiểm các extensions theo README.

Điền host/port/database, tài khoản database và thông tin bootstrap Admin trong môi trường local. Không đưa credentials vào frontend, Git, report hoặc ảnh chụp; tách DB test khỏi dữ liệu ứng dụng.

\subsection{Khởi tạo database và tài khoản}
\begin{lstlisting}[caption={Các bước khởi tạo theo README}]
npm.cmd run backend:connection
npm.cmd run backend:init
npm.cmd run backend:migrate
npm.cmd run backend:create-admin
npm.cmd run backend:seed-store
npm.cmd run backend:doctor
\end{lstlisting}
connection kiểm kết nối, init tạo DB còn thiếu, migrate áp dụng migration mới. create-admin cần BOOTSTRAP\_ADMIN\_EMAIL/PASSWORD trong môi trường, mật khẩu 16–72 byte; không có mật khẩu mặc định hoặc reset Admin cũ. seed-store nạp nội dung công khai một lần, không ghi đè dữ liệu đã sửa.

\subsection{Chạy thủ công với hai terminal}
Terminal thứ nhất chạy PHP:
\begin{lstlisting}
npm.cmd run dev:backend
\end{lstlisting}
Terminal thứ hai chạy frontend:
\begin{lstlisting}
npm.cmd run dev
\end{lstlisting}
Mở \url{http://localhost:8443/}; trang quản trị hiện tại là \url{http://localhost:8443/\#/admin}. PHP mặc định tại \url{http://127.0.0.1:8000}. Vite proxy các đường dẫn API/uploads/products/assets đến PHP. Kiểm server hiện hữu trước khi mở thêm process. Nếu đổi cổng PHP, BACKEND\_PORT và BACKEND\_URL phải tương ứng.

Catalog SSR có thể truy cập trực tiếp tại \url{http://127.0.0.1:8000/products}. Khi backend lỗi, UI phải báo lỗi và cho thử lại; không hiển thị dữ liệu mock như kết quả thật. Kết thúc server bằng Ctrl+C tại terminal tương ứng.

\subsection{Build ứng dụng và triển khai}
\code{npm.cmd run build} tạo frontend/dist; dist riêng không cung cấp PHP API. Bản deploy cần same-origin hoặc reverse proxy đúng các đường dẫn PHP và uploads. DocumentRoot của PHP là backend/public, không phải root repo. Không public .env, .tools, app, config hay vendor; thư mục uploads không được thực thi script.

Production cần HTTPS và cookie Secure, giới hạn upload ở web server/PHP nhất quán và thông báo lỗi không lộ SQL/stack/secret. PHP built-in dev server không được xem là cấu hình production. Backup và việc kiểm phục hồi dữ liệu nên được chuẩn bị trước khi sử dụng dữ liệu thật.

\subsection{Build riêng báo cáo LaTeX}
Chạy từ thư mục report:
\begin{lstlisting}
powershell -NoProfile -ExecutionPolicy Bypass -File .\build.ps1
# Equivalent:
python build.py --tex main.tex --engine lualatex
\end{lstlisting}
PDF: \code{dist/main.pdf}; log: \code{logs/main}.
\clearpage
